
\section{Konkretne przewidywania TGP dla układów wielociałowych}

Niniejsza sekcja zawiera ilościowe wyniki uzyskane za pomocą
dokładnej całki Feynmana (moduł \texttt{three\_body\_force\_exact.py}).
Wszystkie poprzednie obliczenia trójciałowe używały przybliżenia
siodłowego, które ma błąd 160--770\% — wyniki poniżej są pierwszymi
\emph{dokładnymi} przewidywaniami TGP dla tych zagadnień.

%==========================================================================
\subsection{Kiedy efekty trójciałowe dominują? Liczba krytyczna \texorpdfstring{$N_{\rm crit}$}{N\_crit}}
%==========================================================================

\paragraph{Schemat.}
Rozważ $N$ identycznych ciał o sile źródłowej $C$, przy średniej
separacji $d$ (bezwymiarowo $t = m_{\rm sp}\,d$).
Energie wielociałowe skalują się jak:
\begin{align}
V_2^{\rm total} &= \binom{N}{2}\,V_2(d) \sim N^2\,V_2,\\
V_3^{\rm total} &= \binom{N}{3}\,V_3(d,d,d) \sim N^3\,V_3,
\end{align}
zatem stosunek
\begin{equation}
\frac{|V_3^{\rm total}|}{|V_2^{\rm total}|}
= \frac{N-2}{3}\cdot\frac{|V_3^{(1\text{ trojka})}|}{|V_2^{(1\text{ para})}|}
\xrightarrow{N\to\infty} \infty.
\end{equation}
Liczba krytyczna, przy której energia trójciałowa zrównuje się z parową:
\begin{equation}
\boxed{
N_{\rm crit} = \frac{3}{\left|\,V_3^{(1)}/V_2^{(1)}\,\right|} + 2,
}
\end{equation}
gdzie $V_3^{(1)}/V_2^{(1)}$ to stosunek dla jednej trójki/pary, obliczony
dokładnie z całki Feynmana.

\paragraph{Wyniki numeryczne.}
Tabela~\ref{tab:Ncrit} zestawia $|V_3/V_2|$ na jedną trójkę
(trójkąt równoboczny) i odpowiadające $N_{\rm crit}$:

\begin{table}[h]
\centering
\caption{Stosunek $|V_3/V_2|$ (jedna trójka, trójkąt równoboczny)
         i liczba krytyczna $N_{\rm crit}$.
         Parametry: $\beta = \gamma = 1$.}
\label{tab:Ncrit}
\begin{tabular}{ccccc}
\toprule
$t = m_{\rm sp}\,d$ & $C=0.05$ & $C=0.10$ & $C=0.15$ & $C=0.20$ \\
\midrule
1.0 & 1.0\% ($N_{\rm crit}=298$) & 3.5\% (86) & 21\% (\textbf{16}) & 14\% (23) \\
1.5 & 1.6\% (191) & 9.5\% (33) & 14\% (23) & 6.3\% (49) \\
2.0 & 1.8\% (166) & 1.8\% (166) & 1.8\% (166) & 1.8\% (166) \\
2.5 & 0.23\% (1295) & 0.39\% (773) & 0.50\% (599) & 0.59\% (513) \\
3.0 & 0.07\% (4581) & 0.12\% (2499) & 0.17\% (1805) & 0.21\% (1458) \\
\bottomrule
\end{tabular}
\end{table}

\paragraph{Główny wniosek.}
\begin{quote}
Dla $C = 0.15$, $t = m_{\rm sp}\,d = 1.0$ (typowy reżim Yukawa):
$N_{\rm crit} = \mathbf{16}$.
Skupisko zaledwie 16 ciał TGP wykazuje energię trójciałową
porównywalną z parową.
\end{quote}

Jest to przewidywanie \emph{jakościowo różne} od grawitacji Newtonowskiej
(gdzie nie ma trójciałowego członu) i od ogólnej teorii względności
(gdzie trójciałowe poprawki post-Newtonowskie są rzędu $v^4/c^4$,
a nie $C$ jak w TGP).

%==========================================================================
\subsection{Promieniowanie skalarne TGP: odcięcie częstotliwości}
%==========================================================================

\paragraph{Mechanizm.}
Pole TGP jest masywne: $m_{\rm sp} = \sqrt{\beta} > 0$.
Dispersion relation dla fal skalarnych:
\begin{equation}
\omega^2 = k^2 + m_{\rm sp}^2 \geq m_{\rm sp}^2.
\end{equation}
Źródło oscylujące z częstotliwością $\omega < m_{\rm sp}$ nie może
wzbudzić propagujących fal — emituje tylko ewanescentne pole bliskie.
Promieniowanie (eksport energii do nieskończoności) jest możliwe
tylko dla $\omega > m_{\rm sp}$.

\paragraph{Moc promieniowania.}
Dla ciała o przyspieszeniu $a(t) = A\omega^2\cos(\omega t)$ moc
promieniowania skalarnego wynosi:
\begin{equation}
P_{\rm TGP}(\omega)
= \frac{C^2}{6\pi}\,A^2\omega^4\,
\sqrt{1 - \frac{m_{\rm sp}^2}{\omega^2}}\;
\theta(\omega - m_{\rm sp}).
\label{eq:P_TGP}
\end{equation}

\paragraph{Porównanie z grawitacją.}
\begin{itemize}
\item Grawitacja Newtonowska: $P_{\rm GW} \sim G m^2 A^2 \omega^4$ (zawsze).
\item TGP: $P_{\rm TGP} = 0$ dla $\omega < m_{\rm sp}$
  (orbity długookresowe \emph{dokładnie} stabilne).
\item Dla $\omega \gg m_{\rm sp}$: $P_{\rm TGP} \to P_{\rm masless}$
  (jak bezmasowe pole skalarne).
\end{itemize}

\paragraph{Tabela supresji mocy.}

\begin{table}[h]
\centering
\caption{Moc promieniowania TGP \eqref{eq:P_TGP} vs bezmasowe pole,
         dla $C=1$, $A=1$, $m_{\rm sp}=1$.}
\begin{tabular}{rrrr}
\toprule
$\omega$ & $P_{\rm TGP}$ & $P_{\rm bezm.}$ & supresja \\
\midrule
$0.5 < m_{\rm sp}$ & 0 & $3.4\times10^{-3}$ & \textbf{0\%} \\
$0.8 < m_{\rm sp}$ & 0 & $2.1\times10^{-2}$ & \textbf{0\%} \\
$1.0 = m_{\rm sp}$ & $4.4\times10^{-3}$ & $5.4\times10^{-2}$ & 8\% \\
$1.1$ & $3.3\times10^{-2}$ & $7.8\times10^{-2}$ & 42\% \\
$1.5$ & $2.0\times10^{-1}$ & $2.7\times10^{-1}$ & 75\% \\
$2.0$ & $7.3\times10^{-1}$ & $8.4\times10^{-1}$ & 87\% \\
$5.0$ & $3.2\times10^{1}$ & $3.3\times10^{1}$ & 98\% \\
\bottomrule
\end{tabular}
\end{table}

\paragraph{Konsekwencja astrofizyczna.}
Przy $m_{\rm sp} = \sqrt{\beta}$, orbita o okresie $T$ nie promieniuje
jeśli:
\begin{equation}
T > T_{\rm crit} = \frac{2\pi}{m_{\rm sp}} = \frac{2\pi}{\sqrt{\beta}}.
\end{equation}
Dla $\beta = 1$: $T_{\rm crit} \approx 6.28$ (w jednostkach TGP).
Układy o długich okresach orbitalnych (galaktyki, gromady) są
\emph{dokładnie} stabilne w TGP — nie tracą energii na promieniowanie.

%==========================================================================
\subsection{Podsumowanie: trzy nowe przewidywania TGP}
%==========================================================================

\begin{enumerate}

\item \textbf{Próg wielociałowy $N_{\rm crit}$} (wynik dokładny):
  Skupiska $N > N_{\rm crit}$ ciał TGP wykazują energię trójciałową
  porównywalną lub większą od parowej. Dla typowych parametrów
  ($C \sim 0.15$, $t \sim 1$):
  \begin{equation}
  N_{\rm crit} \approx 16 \quad (C=0.15,\; t=1).
  \end{equation}
  \emph{Brak odpowiednika w Newtonie i GR na tym poziomie przybliżenia.}

\item \textbf{Zależność oddziaływania trójciałowego od kształtu}:
  Dwa trójkąty o tym samym obwodzie dają różne $I_Y$.
  Trójkąt zdegenerowany (bardzo skośny) ma $I_Y$ o 19\% większe
  niż równoboczny przy tym samym obwodzie. Przybliżenie siodłowe
  nie widziało tego efektu w ogóle (dawało identyczny wynik).

\item \textbf{Odcięcie promieniowania skalarnego}:
  Orbity z $\omega < m_{\rm sp}$ nie promieniują w TGP.
  Długookresowe układy (gromady, galaktyki) są dokładnie stabilne —
  inaczej niż w GR gdzie każdy układ traci energię przez promieniowanie
  grawitacyjne.

\end{enumerate}

\paragraph{Status obliczeń.}
Wszystkie trzy wyniki uzyskano z dokładnej całki Feynmana
\eqref{eq:IY_feynman_2D} i dokładnego gradientu \eqref{eq:exact_dI_dd},
bez przybliżeń. Poprzednie obliczenia trójciałowe TGP używały wzoru
siodłowego z błędem 160--770\% i brakiem zależności od kształtu.

%==========================================================================
\subsection{Predykcje kosmologiczne TGP-FDM (sektor galaktyczny)}
%==========================================================================

\paragraph{Kontekst.}
Oprócz sektora planckowskiego (Efimov, $m_{\rm sp} \sim 0.1\,l_{\rm Pl}^{-1}$)
TGP zawiera sektor FDM ($m_{\rm sp}^{\rm FDM} = \sqrt{\gamma}
\approx 8.2\times10^{-51}\,l_{\rm Pl}^{-1}$, $m_{\rm boson} \sim 10^{-22}$~eV),
w~którym pole $\Phi$ tworzy solitony ciemnej materii na skalach galaktycznych.

%--------------------------------------------------------------------------
\subsubsection{Kill-shot K18: skalowanie jądra solitonu (ex41--ex44)}
%--------------------------------------------------------------------------

\paragraph{Predykcja F3 (scenariusz uniwersalny).}
Dla stałego $m_{\rm sp} = m_{\rm sp}^{\rm FDM} = \mathrm{const}$ (scenariusz
F3) profil Schive'a~(2014) implikuje:
\begin{equation}\label{eq:k18-prediction}
  \boxed{r_c \propto M_{\rm gal}^{-1/9},
  \quad \alpha_{\rm F3} = -\tfrac{1}{9} \approx -0.111}
\end{equation}
podczas gdy scenariusz F1 ($m_{\rm boson} \propto M_{\rm gal}$) przewiduje
$\alpha_{\rm F1} = -1$.

\paragraph{Wyniki obserwacyjne (ex41--ex44).}
\begin{table}[h]
\centering
\caption{Weryfikacja predykcji K18 ($r_c \propto M_{\rm gal}^{\alpha}$).}
\label{tab:k18}
\begin{tabular}{lcccc}
\toprule
Próba & $n$ & $\alpha_{\rm obs}$ & $\Delta{\rm AIC}({\rm F1{-}F3})$
  & Status F3 \\
\midrule
SPARC (ex41) & 30 & $-0.096\pm0.028$ & $+133$ & OK ($0.5\sigma$) \\
SPARC+THINGS+LTHINGS (ex43) & 75 & $-0.086\pm0.021$ & $+515$
  & \textbf{OK ($1.2\sigma$)} \\
F3+feedback (ex44) & 75 & model M3 & $-294$ vs M0 & \textbf{najlepszy} \\
\bottomrule
\end{tabular}
\end{table}

F1 ($\alpha=-1$) wykluczone na \textbf{44.6\,$\sigma$}.
Napięcie spirale/karłowe 5.2\,$\sigma$ wyjaśnione feedbackiem barionowym:
\begin{equation}
  r_c \propto M_{\rm gal}^{-1/9}\cdot (f_{\rm bar}/f_0)^{\beta_{\rm bar}},
  \quad \beta_{\rm bar} = 0.127, \quad M_{\rm break} = 5.1\times10^9\,M_\odot.
\end{equation}

%--------------------------------------------------------------------------
\subsubsection{Kill-shot K14: napięcie Lyman-$\alpha$ (ex40--ex42)}
%--------------------------------------------------------------------------

\paragraph{Predykcja TGP-FDM.}
Masa bozonu FDM $m_{22} \equiv m_{\rm boson}/(10^{-22}~{\rm eV}) \approx 1$
(z $m_{\rm sp} = \sqrt{\gamma}$, ex39). Predykcja skali supresji widma mocy:
\begin{equation}
  k_{1/2}(m_{22}) = \frac{4.5\,m_{22}^{4/9}}{1\,h/{\rm Mpc}}
  \quad\Rightarrow\quad k_{1/2}(m_{22}=1) \approx 11.3\;h/{\rm Mpc}.
\end{equation}

\paragraph{Napięcie z obserwacjami.}
Irsic+2017 wymagali $m_{22} > 20$ (napięcie 20$\times$).
Rogers+Peiris 2021 (luźniejsze priory $T_{\rm IGM}$): $m_{22} > 2.1$ ---
napięcie \textbf{zredukowane do 2$\times$} (marginalnie).
Zaktualizowana predykcja: $m_{22} = 2.0 \pm 1.5$.

\paragraph{Status.} K14 marginalnie napięty. Rozstrzygnięcie: DESI DR3 (2026+).

\paragraph{Łączona tabela predykcji N-body TGP.}
\begin{table}[h]
\centering
\caption{Pięć falsyfikowalnych predykcji sektora N-body TGP.}
\label{tab:nbody-all}
\begin{tabular}{llll}
\toprule
Predykcja & Wartość & Test & Status \\
\midrule
$N_{\rm crit}$ (efekty 3B) & $\approx 16$ ($C=0.15$, $t=1$) & układy 3-ciałowe & brak danych \\
Kształt a $V_3$ & 19\% dla skośnego vs równobocznego & precyzyjna 3B & brak danych \\
Odcięcie promieniowania & $P_{\rm TGP}=0$ dla $\omega < m_{\rm sp}$ & LIGO/pulsary & OK \\
K18: $\alpha = -1/9$ & $\alpha_{\rm obs} = -0.086\pm0.021$ & SPARC+THINGS & \textbf{OK} (1.2$\sigma$) \\
K14: $m_{22} \approx 2$ & Rogers+2021: $m_{22} > 2.1$ & Lyman-$\alpha$ & marginalnie napięty \\
\bottomrule
\end{tabular}
\end{table}
